\section{Related Work}
\label{sec:related_work}

\subsection{Retrieval-Augmented Generation and Dense Retrieval}
\label{sec:rw_rag_dense}

Retrieval-augmented generation (RAG) grounds language-model outputs in external evidence rather than relying only on parametric memory. The original RAG formulation retrieves passages from a dense vector index and conditions a generator on those passages \cite{lewis2020rag}, while Dense Passage Retrieval (DPR) made dual-encoder retrieval a common backbone for open-domain QA \cite{karpukhin2020dpr}. Later retrieval-augmented readers showed that stronger readers can aggregate information from multiple retrieved passages effectively, as demonstrated by Fusion-in-Decoder \cite{izacard2021fid} and Atlas \cite{izacard2023atlas}. Late-interaction retrievers such as ColBERT improved passage matching beyond single-vector encoders \cite{khattab2020colbert}. RAG surveys organize this space around retrieval quality, augmentation design, and generation faithfulness \cite{gao2023ragsurvey}.

We keep the processed evidence pool, reader prompt, evaluation harness, and top-$k$ setting fixed while comparing multiple retrieval strategies. The canonical MiniLM query-composition variants in Table~\ref{tab:hyde_query_ablation} share the same encoder and FAISS index; each post-hoc BGE diagnostic likewise uses a shared BGE configuration across its compared rows. This positions the work as a controlled multi-hop RAG study focused on query-side expansion.

\subsection{HyDE and Multi-hop Evidence Acquisition}
\label{sec:rw_hyde_multihop}

HyDE is the closest mechanism to our query expansion step. It generates a hypothetical document from the input query and uses that generated text for dense retrieval, without treating the generated document itself as ground-truth evidence \cite{gao2022hyde}. Multi-hop QA benchmarks create a different stress test: HotpotQA emphasizes explainable multi-hop supporting facts \cite{yang2018hotpotqa}, 2WikiMultihopQA evaluates explicit reasoning steps \cite{ho2020twowiki}, and MuSiQue composes multi-hop questions from single-hop relations \cite{trivedi2021musique}. In these settings, the system must recover multiple supporting documents, and one missing hop can prevent the reader from producing the correct answer. Multi-hop dense retrieval learns evidence acquisition over multiple hops \cite{xiong2021multihopdense}, while IRCoT constructs retrieval steps by interleaving chain-of-thought reasoning and retrieval \cite{trivedi2022ircot}.

Recent work has also questioned whether hypothetical-document gains partly reflect benchmark knowledge reproduced by the query generator. Yoon et al. study LLM-based query expansion in fact verification and find that gains are stronger when generated documents contain information entailed by gold evidence, raising benchmark knowledge-leakage concerns \cite{yoon2025knowledgeleakage}. Our distinction from original HyDE and this leakage analysis is the experimental object. We adapt HyDE-style generation to a controlled multi-hop RAG pipeline and test whether richer document-like hypothetical expansion is associated with improved complete supporting-title retrieval. We also evaluate exact answer-string overlap in the generated passage, treating it as a query-side artifact rather than reader evidence.

\subsection{Query Reformulation and Generative Query Expansion}
\label{sec:rw_query_reformulation}

Query reformulation methods improve retrieval by changing the query-side text rather than the retriever architecture. Generation-augmented retrieval expands open-domain QA queries with generated text \cite{mao2020gar}; Query2doc turns a query into an LLM-generated pseudo-document for retrieval \cite{wang2023query2doc}; and Rewrite-Retrieve-Read studies LLM-based query rewriting for retrieval-augmented readers \cite{ma2023rewrite}. ReDE-RF reframes hypothetical-document generation as relevance-feedback selection, explicitly targeting HyDE's dependence on LLM domain knowledge and the cost of generating long hypothetical documents \cite{jedidi2024rederf}. These methods are close to our matched-call control because they also spend an additional query-side generation step before retrieval.

We use these published methods to define the neighboring design space rather than to claim a faithful implementation of each method. The matched-call Single-query Reformulation row is a mechanism control: it asks whether HyDE-style gains can be explained by one extra Qwen3.7-Max call that rewrites the question into a retrieval query. The stricter question-plus-rewritten-query control further asks whether the difference is merely that HyDE retains the original question. Likewise, the equal-budget keyword, direct-rewrite, decomposition, and document-like forms are query-form diagnostics rather than reproductions of Query2doc, Rewrite-Retrieve-Read, or ReDE-RF. In these controlled comparisons, the document-like hypothetical rows obtain higher retrieval and reader scores than direct rewritten-query text. Together with the answer-string-overlap analysis, this diagnostic positions our contribution as a controlled study of fixed-depth, training-free generative expansion for multi-hop evidence acquisition, rather than a new retriever architecture, a full public-baseline benchmark, or an adaptive agent.

\subsection{Agentic and Iterative RAG}
\label{sec:rw_agentic_rag}

Agentic and iterative RAG systems increase flexibility by allowing models to plan, act, revise, or select retrieval strategies over multiple steps. ReAct interleaves reasoning traces with external actions \cite{yao2022react}; IRCoT interleaves chain-of-thought reasoning with retrieval for multi-step questions \cite{trivedi2022ircot}; FLARE actively retrieves when generation confidence indicates missing information \cite{jiang2023flare}; Self-RAG learns reflection signals for retrieval and generation control \cite{asai2023selfrag}; and Adaptive-RAG chooses retrieval strategies according to question complexity \cite{jeong2024adaptiverag}. Recent surveys describe this direction as a shift from static retrieve-then-generate pipelines toward systems with planning and adaptive control \cite{singh2025agenticrag}.

Our pipeline chooses a fixed-depth, training-free design within this broader family. It uses one query-side generated passage, one top-10 retrieval step, and one reader call as the primary retrieval-reader protocol; any guarded answer canonicalization is kept as a supplemental post-processing diagnostic. The rule-based evidence-expanded iterative row in our experiments is therefore a heuristic control for extra retrieval-query executions, not a reproduction of IRCoT's reasoning-and-retrieval procedure. This structure is less expressive than full iterative or agentic RAG, but it makes the cost profile and result ledger easier to inspect.

Post-retrieval evaluators and self-checking methods are adjacent to this adaptive-RAG literature, but they are not a separate method family in our main claim. RAGAS, ARES, and RAGChecker evaluate RAG outputs or diagnose retrieval and generation errors \cite{es2023ragas,saadfalcon2023ares,ru2024ragchecker}; Corrective RAG, RARR, SelfCheckGPT, and Self-RAG use evaluators, retrieved evidence, sampled consistency, or critique signals to improve reliability \cite{yan2024crag,gao2022rarr,manakul2023selfcheckgpt,asai2023selfrag}. These works motivate why answer-format and faithfulness checks are useful after retrieval, but our study focuses on retrieval-stage evidence acquisition. The guarded canonicalization component is therefore only a small exploratory answer-format diagnostic, with its selector, guards, and transition counts reported in the supplement rather than treated as a primary method.
